\section{Distributional losses and recovery bounds}
\subsection{Cross-entropy and the signed shape decomposition}
\begin{proposition}[Strict propriety of CE]\label{prop:ce}
For every true distribution $P$,
\[
 \CE(P,q)=H(P)+\KL(P\Vert q),\qquad \KL(P\Vert q)\ge0,
\]
with equality in the second relation exactly when $q=P$.
Thus expected individual-rating CE uniquely identifies $P$ over the unrestricted simplex.
\end{proposition}
\begin{proof}
Subtracting $H(P)$ from CE gives the displayed KL expression whenever CE is finite; the identity extends to infinite CE.
For $A=\{r:P(r)>0\}$ and a forecast positive on $A$, concavity of the logarithm gives
\[
 -\KL(P\Vert q)=\sum_{r\in A}P(r)\log\frac{q(r)}{P(r)}
 \le\log\sum_{r\in A}q(r)\le0.
\]
Equality requires $q(r)/P(r)$ constant on $A$ and $\sum_{r\in A}q(r)=1$, forcing that constant to be one.
If $q$ is zero on a positive-probability category, KL is infinite and cannot minimize CE.
\end{proof}

\begin{proposition}[Natural exponential families and moment matching]\label{prop:exponential}
Let $h(r)>0$ on $R$ and
\[
 q_\eta(r)=h(r)e^{\eta r-A(\eta)},\qquad
 A(\eta)=\log\sum_r h(r)e^{\eta r},\qquad \eta\in\R.
\]
Then $A'(\eta)=m(\eta)$ and $A''(\eta)=\Var_{q_\eta}(R)>0$.
The mean map is a bijection from $\R$ to $(0,K-1)$.
For any target $P$ with interior mean $t$, expected CE within this family has a unique finite minimizer satisfying $m(\eta)=t$, even if $P$ lies outside the family.
For $h\equiv1$, this mean-matching distribution, denoted $q_t$, uniquely maximizes entropy among all distributions with mean $t$.
\end{proposition}
\begin{proof}
Differentiating the finite sum defining $A$ gives $A'=\sum_r r q_\eta(r)$.
Differentiating again gives $A''=\sum_r r^2q_\eta(r)-(\sum_r rq_\eta(r))^2$.
All categories have positive probability and at least two distinct ratings exist, so this variance is strictly positive.
Hence $m$ is continuous and strictly increasing.
As $\eta\to-\infty$, mass concentrates at zero, and as $\eta\to+\infty$, mass concentrates at $K-1$.
The intermediate value theorem establishes the claimed range and inverse.
Writing $t=\sum_r rP(r)$ gives
\[
 \CE(P,q_\eta)=-\sum_rP(r)\log h(r)-\eta t+A(\eta).
\]
Its first derivative is $m(\eta)-t$ and its second derivative is strictly positive.
The unique zero of the first derivative is therefore the unique global minimizer.
For $h\equiv1$ and $\eta$ chosen to give mean $t$, every distribution $P$ with that mean satisfies
\[
 \KL(P\Vert q_t)=-H(P)-\eta t+A(\eta)=H(q_t)-H(P).
\]
Nonnegativity and the equality condition of KL prove unique entropy maximization.
\end{proof}

At an endpoint target mean, $P$ is the corresponding point mass.
A full-support exponential family approaches that optimum as $\eta\to\pm\infty$; it does not attain it at a finite parameter.
Moment matching here requires the sufficient statistic $r$ and the positive-variance condition.
It is not a property of every one-parameter rating family.
With freely chosen $\eta_i$ per item and interior feasible means, correlation-one predictions remain $m_i=a+bt_i$.
The family's shape is tied to these means, so the remaining ambiguity concerns level and scale.

Let $\eta_t$ and $\eta_m$ have nearby interior means $t$ and $m$, and write $V(t)=A''(\eta_t)>0$.
Then
\begin{align}
 \KL(q_{\eta_t}\Vert q_{\eta_m})
 &=A(\eta_m)-A(\eta_t)-A'(\eta_t)(\eta_m-\eta_t)\\
 &=\frac{(m-t)^2}{2V(t)}+o((m-t)^2).\label{eq:localprice}
\end{align}
Indeed, Taylor expansion gives $\tfrac12 V(t)(\eta_m-\eta_t)^2+o((\eta_m-\eta_t)^2)$, and the inverse-function derivative gives $\eta_m-\eta_t=(m-t)/V(t)+o(m-t)$.
For a misspecified target of mean $t$, subtracting its CE at $\eta_t$ from its CE at $\eta_m$ gives the same exact expression on the first line.
This is a local cost of mean miscalibration inside this family.
It does not assert that the family's variance equals the target variance.

\begin{proposition}[Signed CE decomposition]\label{prop:cross}
Let $P$ have interior mean $t$ and let $q_t$ be its maximum-entropy mean-matching distribution from Proposition~\ref{prop:exponential} with $h\equiv1$.
For a forecast $q$ positive on every category, define
\[
 G(P)=H(q_t)-H(P)\ge0,\qquad
 C(P,q)=\sum_r(P(r)-q_t(r))(-\log q(r)).
\]
Then
\begin{equation}\label{eq:cross}
 \CE(P,q)=H(P)+G(P)+\KL(q_t\Vert q)+C(P,q).
\end{equation}
The cross term $C(P,q)$ can be negative.
If $\log q(r)=\alpha+\beta r$, it vanishes, and the minimum excess CE within that log-affine-in-rating family is $G(P)$.
\end{proposition}
\begin{proof}
Add and subtract $\CE(q_t,q)$ to obtain $\CE(P,q)=\CE(q_t,q)+C(P,q)$.
Proposition~\ref{prop:ce} gives $\CE(q_t,q)=H(q_t)+\KL(q_t\Vert q)$.
Substitution of $H(q_t)=H(P)+G(P)$ yields Eq.~\eqref{eq:cross}.
If $\log q(r)=\alpha+\beta r$, then
\[
 C(P,q)=-\alpha\sum_r(P(r)-q_t(r))-\beta\sum_r r(P(r)-q_t(r))=0
\]
because both total mass and mean agree.
KL is uniquely minimized at $q=q_t$, establishing the family-specific excess $G(P)$.
For any positive $P$, substituting the unrestricted forecast $q=P$ into Eq.~\eqref{eq:cross} yields
\[
 C(P,P)=-G(P)-\KL(q_t\Vert P).
\]
This is negative whenever $P\ne q_t$.
\end{proof}

The nonnegative shape gap is therefore a restriction cost for the stated log-affine family, not an unavoidable penalty for every forecast.
The cross term prevents interpreting Eq.~\eqref{eq:cross} as a sum of nonnegative errors.
Other scalar-to-distribution families require their own analysis.
At endpoint means, the target and its maximum-entropy representative coincide and $G(P)=C(P,q)=0$.

\subsection{Ranked probability score and distribution distances}
\begin{proposition}[Expected RPS, signed CDF identity, and decomposition]\label{prop:rps}
For any $P,q$, let $d_k=F_P(k)-F_q(k)$, $\bar d=(K-1)^{-1}\sum_k d_k$, and $D(P,q)=\sum_k d_k^2$.
Then
\begin{align}
 \E_P S(q,Y)&=\frac1{K-1}\sum_k F_P(k)(1-F_P(k))+\frac{D(P,q)}{K-1},\label{eq:riskrps}\\
 \sum_{k=0}^{K-2}d_k&=m_q-m_P,\label{eq:signedcdf}\\
 \mathcal E_P(q)&=\frac{(m_q-m_P)^2}{(K-1)^2}+\frac1{K-1}\sum_k(d_k-\bar d)^2.\label{eq:rpssplit}
\end{align}
Consequently, individual-rating normalized RPS is strictly proper.
Also $|m_q-m_P|\le(K-1)\TV(P,q)$.
\end{proposition}
\begin{proof}
For $Z_k=\mathbf1\{Y\le k\}$, $\E Z_k=F_P(k)$ and $\Var(Z_k)=F_P(k)(1-F_P(k))$.
The identity $\E(a-Z_k)^2=(a-\E Z_k)^2+\Var(Z_k)$, summed and divided by $K-1$, proves Eq.~\eqref{eq:riskrps}.
Interchanging the finite sums gives
\[
 \sum_{k=0}^{K-2}F_P(k)=\sum_rP(r)(K-1-r)=K-1-m_P.
\]
Subtracting its counterpart for $q$ proves Eq.~\eqref{eq:signedcdf}, including its sign.
The centered sum-of-squares identity is
\[
 \sum_kd_k^2=\frac{(\sum_kd_k)^2}{K-1}+\sum_k(d_k-\bar d)^2.
\]
Dividing by $K-1$ gives Eq.~\eqref{eq:rpssplit}.
The excess vanishes exactly when all CDF values agree, which determines every category probability by successive differences.
This proves strict propriety, including distributions on the simplex boundary.
Finally, mass conservation gives
\[
 |m_q-m_P|=\left|\sum_r\left(r-\frac{K-1}{2}\right)(q(r)-P(r))\right|
 \le\frac{K-1}{2}\sum_r|q(r)-P(r)|.
\]
\end{proof}

The first term of Eq.~\eqref{eq:riskrps} is the irreducible expected score for a fresh stochastic rating.
It is distinct from excess RPS, which can be zero.
Equation~\eqref{eq:rpssplit} has mean-error denominator $(K-1)^2$, which is $25$ on a six-category scale.
The remaining term measures variation of the CDF differences after their average is removed.
It vanishes exactly when all CDF differences are equal; it is not an independent variance error.
If means agree, all remaining CDF discrepancy appears in this term.

TVD to the true distribution equals zero exactly when $q=P$.
It is a distribution distance rather than a strictly proper individual-rating score.
In fact,
\[
 \TV(q,\delta_y)=\frac12\left(1-q(y)+\sum_{r\ne y}q(r)\right)=1-q(y),
\]
so $\E_{Y\sim P}\TV(q,\delta_Y)=1-\sum_rP(r)q(r)$.
This expected loss is minimized by forecasts supported on the modal categories of $P$.
It generally does not recover $P$.
TVD to a sparse empirical histogram is also a noisy diagnostic, whose expectation need not equal TVD to $P$.

\begin{proposition}[CE controls bounded moments]\label{prop:pinsker}
For every $f:R\to[0,1]$,
\begin{equation}\label{eq:momentbound}
 |\E_P f-\E_q f|\le\TV(P,q)\le\sqrt{\KL(P\Vert q)/2}.
\end{equation}
In particular, the bound applies to every normalized raw moment $f(r)=(r/(K-1))^j$, $j\ge1$.
\end{proposition}
\begin{proof}
Write $\delta(r)=P(r)-q(r)$ and $A=\{r:\delta(r)\ge0\}$.
Since $\sum_r\delta(r)=0$, $\sum_{r\in A}\delta(r)=\TV(P,q)=d$.
The positive and negative parts bound $\sum_r f(r)\delta(r)$ above by $d$ and below by $-d$.
To prove the second inequality, suppose KL is finite; otherwise it is immediate.
Grouping categories in $A$ and its complement cannot increase KL.
For completeness, for any group $B$, concavity of the logarithm applied with weights $P(r)/P(B)$ yields
\[
 \sum_{r\in B}P(r)\log\frac{P(r)}{q(r)}\ge P(B)\log\frac{P(B)}{q(B)}.
\]
Zero group masses follow by continuity.
Set $a=P(A)$ and $b=q(A)$, so $a-b=d$.
The grouped bound is the Bernoulli KL $h_b(a)$.
For $a,b\in(0,1)$, $h_b(b)=h_b'(b)=0$ and
\[
 h_b''(a)=\frac1{a(1-a)}\ge4.
\]
Integrating twice between $b$ and $a$ gives $h_b(a)\ge2(a-b)^2=2d^2$.
Endpoint cases follow by limits or have infinite KL.
Thus $\KL(P\Vert q)\ge2\TV(P,q)^2$.
\end{proof}

\subsection{Limits of recovering distributions from means alone}
\begin{proposition}[Two-distribution recovery lower bounds]\label{prop:lower}
Suppose a recovery rule receives only the exact mean and outputs a distribution $Q(m)$.
Let $P_0,P_1$ be distinct distributions with the same mean, and set $d=\TV(P_0,P_1)$.
For $q=Q(m)$,
\begin{align}
 \max_{j\in\{0,1\}}\TV(P_j,q)&\ge d/2,\\
 \frac{\mathcal E_{P_0}(q)+\mathcal E_{P_1}(q)}2
 &\ge\frac{D(P_0,P_1)}{4(K-1)},\label{eq:twopointrps}\\
 \max_{j\in\{0,1\}}\KL(P_j\Vert q)&\ge d^2/2.
\end{align}
The first two bounds are sharp when the feasible forecast set includes $(P_0+P_1)/2$.
\end{proposition}
\begin{proof}
The rule sees the same input under both truths and therefore must use the same forecast.
The triangle inequality gives $d\le\TV(P_0,q)+\TV(q,P_1)$, proving the first bound.
Let $v_j=(F_{P_j}(k))_{k=0}^{K-2}$ and $u=(F_q(k))_{k=0}^{K-2}$.
The Euclidean identity
\[
 \frac{\|v_0-u\|^2+\|v_1-u\|^2}{2}
 =\left\|u-\frac{v_0+v_1}{2}\right\|^2+\frac{\|v_0-v_1\|^2}{4}
\]
and Proposition~\ref{prop:rps} give Eq.~\eqref{eq:twopointrps}.
The mixture forecast has the midpoint CDF and TVD $d/2$ to each truth, proving sharpness of those bounds.
For at least one truth the TVD is at least $d/2$; Proposition~\ref{prop:pinsker} then gives KL at least $2(d/2)^2$.
\end{proof}

For a randomized rule, the RPS bound holds after averaging over its randomization, and the corresponding TVD and KL worst-case bounds hold for expected losses.
Indeed the same inequalities can first be averaged over the common output law, using $\max_j\E L_j\ge(\E L_0+\E L_1)/2$ and convexity where necessary.
For raw expected RPS, the right side of Eq.~\eqref{eq:twopointrps} must be added to the average of the two irreducible scores in Eq.~\eqref{eq:riskrps}.
These are information bounds for a rule given only means.
They do not lower-bound an estimator that also receives informative individual ratings, item features, or valid shape restrictions.
Nor do they force every correlation optimizer to choose an inaccurate member of its optimum set.
